\chapter*{Conclusion Générale et Perspectives}
\addcontentsline{toc}{chapter}{Conclusion Générale et Perspectives}

Ce projet de fin d'études nous a permis de mener à bien un cycle complet de
développement logiciel : de l'analyse du besoin métier jusqu'au déploiement d'une
solution opérationnelle sur une infrastructure d'entreprise réelle.

En partant d'un constat simple — la gestion des incidents informatiques au sein
du Groupe Informatique de la Branche Carburants de Naftal souffrait d'un manque
criant de traçabilité, d'organisation et de visibilité — nous avons conçu et réalisé
une \textbf{Plateforme de Gestion des Incidents Applicatifs (GIA)} qui répond
intégralement au cahier des charges validé par la Direction Informatique.

La solution s'appuie sur une architecture \textbf{3-Tiers} éprouvée, combinant
une interface Bootstrap~5 moderne et responsive, une logique métier PHP~8
sécurisée par un système de rôles et de sessions, et une base de données relationnelle
SQL Server garantissant l'intégrité et la traçabilité complète via une table d'audit
(\texttt{incident\_logs}). Le déploiement sur \textbf{Windows Server~2019 / IIS}
confirme la conformité de la solution avec les standards d'infrastructure de Naftal.

Les principaux apports de ce projet sont :
\begin{itemize}
    \item \textbf{Traçabilité complète :} chaque action sur un ticket est horodatée
    et attribuée à un utilisateur identifié, rendant le système opposable lors d'un audit.
    \item \textbf{Pilotage par les données :} les KPIs du tableau de bord permettent
    aux managers de mesurer objectivement la performance du support (MTTR, taux de
    résolution, charge par technicien).
    \item \textbf{Workflow structuré :} la machine à états du cycle de vie des tickets
    élimine les risques d'incidents oubliés ou traités en double.
\end{itemize}

\section*{Limites du travail et positionnement par rapport au MEF}
Ce mémoire n'applique pas le modèle d'évaluation fonctionnelle (MEF)~\cite{autissier2008} sous
la forme d'enquêtes standardisées auprès de la DCSI : les baromètres (activités, compétences,
organisation, satisfaction) n'ont pas été mesurés dans le cadre du PFE. En revanche, le cadre
théorique et les travaux sur Naftal~\cite{meghari2014} ont guidé la \textbf{justification} du
projet : la rubrique «\,maintenance applicative\,» et la «\,gestion de la relation
utilisateurs\,» appellent des dispositifs traçables. GIA constitue une brique concrète ; une
évaluation complète de la performance du support à l'échelle groupe exigerait des indicateurs
consolidés, des entretiens élargis et une gouvernance des données sur la durée.

\paragraph*{Limites} Cette première version de GIA présente néanmoins certaines limites qu'il
convient d'énoncer honnêtement. Premièrement, la solution a été validée sur un environnement
simulé (VirtualBox) et non sur le réseau de production réel de Naftal : des tests de charge et de
montée en charge restent nécessaires avant tout déploiement à grande échelle. Deuxièmement, la
plateforme ne couvre actuellement que les incidents applicatifs ; les incidents matériels (panne
d'imprimante, défaillance réseau) en sont exclus. Troisièmement, l'absence de module de
notification automatique contraint les utilisateurs à consulter activement la plateforme, limitant
la réactivité.

\section*{Apports pour le parcours académique}
Sur le plan méthodologique, le projet illustre le passage du constat qualitatif (entretiens,
cahier des charges, observation) à une spécification testable (règles R1--R7, matrices de
traçabilité), puis à une réalisation contrôlée sur stack imposée (Windows Server, IIS, PHP, SQL
Server). Cette chaîne reproduit les étapes rencontrées dans les projets SI en entreprise et
complète la vision «\,macro\,» de la fonction support développée dans les chapitres théoriques.

\section*{Perspectives}
La plateforme GIA constitue une base solide qui peut être enrichie dans le cadre de la stratégie
de modernisation du SI de Naftal.

\paragraph*{Notifications automatiques}
L'intégration d'une couche de notification par e-mail via le serveur Exchange de Naftal
permettrait d'alerter en temps réel le Demandeur et l'Administrateur à chaque transition de
statut critique. Cela réduirait le délai de réaction et éliminerait le besoin de consulter
manuellement l'interface pour suivre l'état d'un ticket.

\paragraph*{Base de connaissances}
À terme, les rapports de résolution rédigés par les techniciens constituent une matière première
précieuse. Une fonctionnalité de recherche plein texte sur les clôtures passées permettrait de
créer une FAQ opérationnelle, réduisant le MTTR sur les incidents récurrents.

\paragraph*{Extension multi-branches}
L'architecture intranet VPN de Naftal (2 Mb/s, reliant toutes les branches) est déjà en place. La
généralisation de GIA aux Branches GPL et Commercialisation ne nécessiterait qu'une adaptation des
catégories d'incidents et des droits d'accès, sans refonte architecturale.
